\chapter{Listagens Representativas dos Módulos de Simulação}
\label{ap:listagens}

\paragraph{}Este apêndice reproduz listagens \emph{representativas} da implementação dos módulos de simulação do \textit{MinervaMesh}. O código-fonte integral de todos os módulos --- incluindo as interfaces HTML e os arquivos de configuração de dependências --- está disponível no repositório público da plataforma, sob licença MIT, em \url{https://github.com/lgsfarias/minervamesh}. A Tabela~\ref{tab:repo_modulos} mapeia cada módulo ao seu diretório no repositório; em cada diretório, a lógica numérica reside no arquivo \texttt{simulation.py}. As listagens selecionadas a seguir cobrem os dois padrões de implementação da plataforma: um módulo 1D completo (condução permanente, o caso mais didático) e o laço temporal do solver de Navier-Stokes (o caso mais complexo), este último precedido de uma descrição do núcleo numérico compartilhado pelos módulos 2D.

\begin{table}[H]
\centering
\caption{Localização de cada módulo de simulação no repositório do \textit{MinervaMesh}. Caminhos relativos ao diretório \texttt{simulations/} na raiz do repositório.}
\label{tab:repo_modulos}
\footnotesize
\begin{tabular}{ll}
\hline
\textbf{Módulo} & \textbf{Diretório (sob \texttt{simulations/})} \\
\hline
Condução permanente 1D (MEF)     & \texttt{mef/conducao-permanente-barra1d-mef/} \\
Condução transiente 1D           & \texttt{mef/calor-transiente-1d/} \\
Convecção-difusão 1D (SUPG)      & \texttt{mef/conveccao-difusao-1d/} \\
Viga 1D (Euler-Bernoulli)        & \texttt{mef/viga-1d-mef/} \\
Vibração 1D (autovalor)          & \texttt{mef/vibracao-1d-mef/} \\
Transferência de calor 2D        & \texttt{mef/transcalor/} \\
Escoamento potencial 2D          & \texttt{mef/escoamento-fluido-2d/escoamento-potencial/} \\
Navier-Stokes 2D                 & \texttt{mef/escoamento-fluido-2d/navier-stokes/} \\
Núcleo MEF 2D compartilhado      & \texttt{mef/escoamento-fluido-2d/common.py} e \texttt{geometry.py} \\
Condução permanente 1D (MDF)     & \texttt{mdf/conducao-permanente-barra1d-geracao/} \\
Equação da onda 1D (MDF)         & \texttt{mdf/equacao-onda/} \\
\hline
\end{tabular}
\end{table}

\section{Módulo 1D Completo: Condução Permanente}
\label{ap:lst_conducao1d}

\paragraph{}A listagem a seguir reproduz integralmente o \texttt{simulation.py} do módulo de condução permanente 1D, representativo da estrutura comum a todos os módulos unidimensionais: leitura de parâmetros do DOM, montagem das matrizes elementares e globais, imposição de condições de contorno, solução do sistema e renderização do resultado no navegador.

\lstinputlisting[
  caption={Condução permanente 1D --- \texttt{simulation.py} (completo)},
  label={lst:conducao1d}
]{../simulations/mef/conducao-permanente-barra1d-mef/simulation.py}

\section{Núcleo Numérico dos Módulos 2D}
\label{ap:lst_nucleo2d}

\paragraph{}O núcleo numérico compartilhado pelos solvers de escoamento 2D concentra-se no módulo \texttt{common.py}, acompanhado de \texttt{geometry.py}. Esse núcleo é responsável pela geração da malha triangular por triangulação de Delaunay e pela montagem das matrizes globais do MEF --- rigidez, massa e advecção --- a partir das contribuições elementares. Por ser código de infraestrutura, ele é reutilizado integralmente tanto pelo módulo de escoamento potencial quanto pelo de Navier-Stokes, evitando duplicação da lógica de montagem entre os dois solvers. Optou-se por não reproduzir aqui a listagem completa desse núcleo, dada a sua extensão; o código integral está disponível no repositório público, na pasta \texttt{simulations/mef/escoamento-fluido-2d/}, em \url{https://github.com/lgsfarias/minervamesh}.

\paragraph{}Do solver de Navier-Stokes 2D, por outro lado, reproduz-se o trecho central: o laço de avanço temporal da formulação Vorticidade-Corrente, com a solução da equação de Poisson para $\psi$, a imposição da condição de contorno de Thom nas paredes, a recuperação do campo de velocidades e o transporte semi-implícito da vorticidade.
% Trecho extraído de simulations/mef/escoamento-fluido-2d/navier-stokes/simulation.py,
% função solve_navier_stokes, laço "while current_step < total_steps" (linhas 209-260).
% Ao editar os comentários do simulation.py, conferir e reajustar o linerange abaixo.

\lstinputlisting[
  caption={Navier-Stokes 2D --- laço temporal da formulação Vorticidade-Corrente (trecho de \texttt{simulation.py})},
  label={lst:ns2d_loop},
  linerange={209-260},
  firstnumber=209
]{../simulations/mef/escoamento-fluido-2d/navier-stokes/simulation.py}
